\chapter{Introduction}

\section{Background and Motivation}
Financial difficulties often develop over several years before a company goes bankrupt. Changes in profitability, debt, and liquidity may offer signs of distress, making bankruptcy prediction a useful application of machine learning. As industrial economics (INDØK) students with an interest in economics and finance, we chose a corporate bankruptcy dataset to connect these interests with the methods taught in IKT466, ``Introduction to Machine Learning.''

A company's latest financial figures provide a snapshot of its situation, while earlier observations show how that situation has changed. Historical data may help distinguish a temporary setback from a longer decline. However, more data and a more complex model do not necessarily produce better predictions. We therefore compare several modelling approaches using both the latest observation and the full financial history available for each company.

\section{Objective and Research Questions}
The objective is to examine whether financial history improves bankruptcy prediction and whether its usefulness differs across model families. We compare logistic regression, XGBoost, a multilayer perceptron (MLP), and a Gated Recurrent Unit (GRU) network. Each model is trained separately under two input conditions: the latest recorded year and all available recorded years ending at that same year. This gives eight model--input combinations.

Our main research question is:
\begin{quote}
Does using all available financial history improve corporate bankruptcy prediction compared with using only the latest year, and how does this effect differ across model families?
\end{quote}

We address this through two investigations:
\begin{enumerate}
    \item Compare the four model families within each input condition, so that they receive the same financial information.
    \item Compare latest-year and full-history performance within each model family, including an analysis by the amount of history available.
\end{enumerate}

The GRU processes annual observations as a sequence. Logistic regression, XGBoost, and MLP receive the same observations arranged as fixed-width feature vectors. These experiments assess the benefit of using the available history; they do not identify an optimal number of historical years or assume that the GRU will perform best.

\section{Scope and Approach}
We use the US public companies bankruptcy dataset, which contains annual financial observations and company-level \texttt{alive}/\texttt{failed} labels. Each company contributes one outcome anchored at its last recorded year. Under the publisher's convention, a failed company's terminal financial year precedes its bankruptcy filing. Filing dates and complete follow-up for alive companies cannot be independently verified from the CSV, so this interpretation remains a limitation of the study.

We use one fixed company-level split into training, validation, and test sets. The split is stratified jointly by outcome and recorded-history group to maintain similar proportions of failed companies and short and long histories across partitions. The same companies, anchor years, and outcomes are used for all eight combinations. Companies with short histories are retained. Missing calendar years and padding are represented explicitly; having fewer records is not assumed to mean that a company is newly established.

Optuna is used to select hyperparameters and class-weight settings separately for each model--input combination by maximising validation Average Precision. We use the same search method and trial budget, with model-specific search spaces that are kept consistent between input conditions where applicable. The selected settings may therefore differ between latest-year and full-history models. Model coefficients and neural network weights are learned from the training data within each trial; Optuna selects the settings used for that training. Preprocessing is fitted on training data only, and classification thresholds are selected on validation data before final test evaluation.

Because bankrupt companies form a small share of the dataset, we use Average Precision as the main metric, supported by precision, recall, and F1. We examine missed bankruptcies, false alarms, training behaviour, and computing cost to interpret the results. Comparisons by history group are reported with their sample sizes, since some groups contain few failed companies. The study evaluates unseen companies within this dataset and does not establish performance in a later economic period or the effect of history independently of changes in tuning, data availability, and model capacity.
